\documentclass[11pt]{article}
% Built with tectonic (XeTeX): cd docs/problems && tectonic -X compile article.tex
% Same page design as the Lean article (docs/lean/article.tex): Charter, 27 mm margins, 1.1 spacing.
\usepackage{fontspec}
\IfFontExistsTF{Charter}{\setmainfont{Charter}}{\setmainfont{XCharter}}
\IfFontExistsTF{Menlo}{\setmonofont{Menlo}[Scale=0.82]}{\setmonofont{DejaVu Sans Mono}[Scale=0.82]}
\usepackage{polyglossia}\setmainlanguage{english}\setotherlanguage{french}
\usepackage{microtype}
\usepackage[a4paper,margin=27mm]{geometry}
\usepackage{setspace}\setstretch{1.1}
\usepackage{amsmath,amssymb,amsthm}
\usepackage[most]{tcolorbox}
\usepackage{titlesec}
\titlespacing*{\section}{0pt}{2.2ex plus 1ex}{1.2ex}
\usepackage{xcolor}
\usepackage[hidelinks]{hyperref}
\usepackage{xurl}
\usepackage{enumitem}
\setlist{nosep,leftmargin=*}
\setlength{\emergencystretch}{2em}
\setlength{\parskip}{0.35em}
\newcommand{\q}[1]{«\,#1\,»}
\newcommand{\fac}[1]{\href{https://grothendieck.umontpellier.fr/#1.pdf}{folder~#1}}
\newcommand{\cd}{\operatorname{cd}}
\newcommand{\Ext}{\operatorname{Ext}}
\newcommand{\Pic}{\operatorname{Pic}}
\newcommand{\Spec}{\operatorname{Spec}}
\newcommand{\Top}{\operatorname{Top}}
\newcommand{\Ind}{\operatorname{Ind}}
\newcommand{\Gal}{\operatorname{Gal}}
\theoremstyle{definition}
\newtheorem{problem}{Problem}
\theoremstyle{plain}\newtheorem*{proposition}{Proposition}\theoremstyle{definition}
% Each problem statement in a frame, as the statements of the Lean article.
\tcolorboxenvironment{problem}{enhanced, breakable, colback=white, colframe=black!55, boxrule=0.6pt, arc=0pt, left=7pt, right=7pt, top=4pt, bottom=4pt, before skip=14pt, after skip=8pt}
% The large conjectures of \S\ref{sec:big}, in a lighter frame.
\newtheorem{bigconj}{Conjecture}
\tcolorboxenvironment{bigconj}{enhanced, breakable, colback=black!3, colframe=black!35, boxrule=0.4pt, arc=0pt, left=7pt, right=7pt, top=4pt, bottom=4pt, before skip=14pt, after skip=8pt}
\newcommand{\entry}[1]{\par\medskip\noindent\textbf{#1}\enspace}
\newcommand{\src}[1]{\entry{Source.}#1}
\newcommand{\ctx}{\entry{Context.}}
\newcommand{\known}{\entry{What is known.}}
\newcommand{\ideas}{\entry{Ideas.}}
\newcommand{\leanst}{\entry{Lean.}}
\newcommand{\settle}{\entry{What would settle it.}}

\title{Open questions from Grothendieck's Montpellier notes:\\ fourteen statements a formalisation could not reach}
\author{Michel Hua\\ \small Commutator, Paris\\ \small\texttt{michel@commutator.io} · ORCID 0009-0002-8304-0683}
\date{Working paper, 29 October 2026}

\begin{document}



\maketitle

\begin{abstract}
The University of Montpellier has put online 178 folders of Alexandre Grothendieck's mathematical working notes. Since August 2026 an edition has transcribed 165 of them with a language model, restated each in current notation, searched the literature for the statements they contain, and formalised what it could in Lean~4 against Mathlib. Most of the statements that look new sit in theories Mathlib does not have yet: étale cohomology, topoi, stacks, abelian schemes, motives. This note collects fourteen of them for mathematicians. Each comes with its folder and page, the context, what the literature search found, the argument the page gives or sketches, the steps the edition added, and what would settle it. A last section takes seven conjectures that are in print under his name, from the standard conjectures to the $p$-curvature conjecture, and drafts from the pages the route he meant to take to each and where it stops. None of the statements has been checked by a specialist, and several may already be in print; a reference is as welcome as a proof or a counterexample.
\end{abstract}


\section{Why this note}

Grothendieck gave his mathematical papers to Jean Malgoire in 1990 and 1995. They were deposited at the University of Montpellier in 2010 and given to it in 2012, and the university placed 178 folders, 16,074 pages by its inventory's count, in open access as facsimiles. They are working notes: drafts in ink and pencil, restarted, struck, annotated, almost never dated. Twenty folders have been edited by hand. The rest had not been transcribed.

Since August 2026 an edition has read the remaining folders with a language model as first transcriber, and written for each folder a \emph{modernised reading}, a restatement in current notation made from the transcription and never from the hand. A register of \emph{findings} lists the statements of the folders that may not stand in the literature, each written as a claim about the literature with what was searched and what would settle it; at the time of writing it holds 456 rows, of which 83 turned out to be in print. Statements were then formalised in Lean~4 against a pinned revision of Mathlib, about two thousand theorems in all. A companion paper describes the formalisation~\cite{huacicm}.

Lean settles a statement only when the objects it speaks of exist in the library. In October 2026, 203 open findings with no proof were sorted by what Mathlib can reach: 16 were reachable as it stands, 39 in part and 147 not at all, and one was set aside as ill-posed; the 147 are out of reach because they live in étale or crystalline cohomology, topoi, stacks, derivators, motives or Hodge theory~\cite{huacicm}. Most of Grothendieck's own work was in those theories. This note hands the most promising of them to people who can prove them with pen and paper, or say where they are already proved.

\subsection*{How to read the entries}

Each problem gives its \emph{source}: the folder (a link to Montpellier's facsimile), the pages in the archive's pagination, and the inventory's dating, whose brackets mark an inferred date. \q{The page} means what Grothendieck wrote, as transcribed; short quotations are his. \q{The reading} means the edition's modernised reading, made by a language model. \q{We} marks a step of the edition: an argument, a correction or a search made for this note or for the register, with a language model, and not checked by a specialist unless the entry says so. The entries make no claim of priority. A statement on an undated leaf says when he wrote it down and nothing about whether others had.

The transcriptions themselves are not public: the rights belong to Grothendieck's heirs, and the reading views stay off the site until they authorise them. The facsimiles are public, at Montpellier, and each entry cites the page to look at. Montpellier's servers were down at the end of October 2026; the links follow the archive's own scheme and should work again.

\section{Étale cohomology and Hodge theory}

\begin{problem}[A global bound with a shift of 3]\label{p:cd}
For a field $k$ of characteristic $p$ put $p^{\nu(k)} = [k : k^{p}]$, and for a scheme $Z$ of characteristic $p$ let $\nu(Z)$ be the supremum of $\nu(k(z))$ over its maximal points. Let $X$ be an affine scheme on which $p$ is not invertible, so that $V(p) \subset X$ has characteristic $p$. Is it true that
\[
\cd_p(X) \;\le\; \sup\Bigl(\sup_{x \in X} \cd_p k(x),\ \nu(V(p)) + 3\Bigr)\,?
\]
\end{problem}

\src{\fac{31}, p.~84, \emph{SGA A [SGA 4] (résidus de rédaction)}, inventory dating [1963--1973]. The page is headed \q{Conjectures sur la cohomologie des schémas affines}; this is its conjecture~2°). The page writes $\ell$ for the prime; we write $p$, since it is the characteristic of $V(p)$ and not a prime invertible on $X$.}

\ctx M.~Artin wrote in SGA~4 X~3.0 that in unequal characteristic one must \q{certainement tenir compte} of the rank of $\Omega^1_k$, without a bound \cite{sga4}. The same page states a local form, 4°): for $U$ open in the spectrum of a henselian local domain $A$ of residue characteristic $p$ with fraction field of characteristic~$0$, $\cd_p(U) \le \dim A + \nu(k) + \cd_p(k)$. The words around \q{ouvert} in 4°) are hard to read, so whether $U$ is meant to be affine, and whether it may meet $V(p)$, is not known.

\known The local form holds, in a sharper form, for $U$ a non-empty affine open of $\Spec A[1/p]$ with $A$ an excellent henselian local domain and $\operatorname{char} K = 0$: Gabber and Orgogozo prove $\cd_p(U) = \dim A + \dim_p(k)$ \cite[Th.~6.1]{gabberorgogozo}, and Kato's $p$-dimension satisfies $\dim_p(k) \le \nu(k) + \cd_p(k)$, which is immediate from their Définition~1.1. Mathew proves that the residue fields of $\Spec R[1/p]$ have $\cd_p \le d+1$ when $R$ is an excellent normal domain, henselian along an ideal $I \ni p$, with $\dim R_{\mathfrak p} + \log_p [\kappa(\mathfrak p) : \kappa(\mathfrak p)^p] \le d$ for $\mathfrak p \supseteq I$; the same bound for the affine opens $\Spec R[1/f]$ of $\Spec R[1/p]$ is established in the course of the proof \cite[Prop.~4.5]{mathew}. Globally, SGA~4 gives $\cd_p X \le 2\dim X + 1$ for $X$ of finite type over $\mathbb Z$ \cite[X~6.2]{sga4}, and Clausen and Mathew give $\cd_p(X) \le \dim X + \sup_x \cd_p k(x)$ for $X$ qcqs of finite Krull dimension \cite[Cor.~3.29]{clausenmathew}; for $\Spec \mathbb Z_p[x_1,\dots,x_n]$, which is not of finite type over $\mathbb Z$, the latter gives $2n+3$ where the page's bound gives $n+3$. Gabber and Zavyalov settle the analogous expectation for affinoids \cite[Th.~1.7]{gabberzavyalov}; the Travaux de Gabber treat $\ell$ invertible only.

\ideas The constant cannot be lowered to~$2$ (our observation): for $X = \Spec \mathbb Z$ and $p$ odd, $\nu(V(p)) = 0$ and $\sup_x \cd_p k(x) = \cd_p \mathbb Q = 2$, while $\cd_p(\Spec \mathbb Z) = 3$ \cite[X~6.1]{sga4}; by Artin--Verdier duality $H^3(\Spec \mathbb Z, j_!\mu_p) \cong \mathbb Z/p$ \cite[II.3.2, 3.3]{milneadt}. A natural route, which is the one of Mathew's proof, is the excision sequence for $V(p) \subset X$ and its complement $X[1/p]$: for $X$ noetherian, $V(p)$ is affine of characteristic~$p$, so of $p$-cohomological dimension at most~$1$ \cite[X~5.1]{sga4} (for $X$ not noetherian a limit argument is needed); local cohomology along $V(p)$ can be moved to the henselisation along $V(p)$ by Gabber's affine analogue of proper base change, where Gabber--Orgogozo and Mathew apply. The difficulty is to absorb the dimension terms of those local bounds into $\sup_x \cd_p k(x)$, which the page's bound allows, and to see where the shift~$3$ comes from: the page does not explain it.

\leanst Not attempted: Mathlib has no étale cohomology.

\settle A printed bound of this shape for an arbitrary affine scheme of mixed characteristic, or a counterexample to the~$3$.

\begin{problem}[Lower bound for a local domain: the exceptions]\label{p:lower}
Let $A$ be a noetherian local domain with fraction field $K$ and residue field $k$, and $r$ a prime. The page states that
\[
\cd_r(K) \;\ge\; \cd_r(k) + \dim A
\]
with exactly two possible exceptions: (a)~$r$ is the characteristic of $K$; (b)~$r = 2$ and some point $x$ of $\Spec A$ has an orderable residue field $k(x)$ with $\cd_2 k(x)(\sqrt{-1})$ finite. It adds that exception~(b) disappears when $A$ is regular, and reduces to \q{$k$ itself real} when $A$ is catenary and universally japanese. Are these refinements true?
\end{problem}

\src{\fac{31}, p.~22, inventory dating [1963--1973]. The margin gives the example that keeps exception~(b) alive: the local ring at the origin of $\mathbb R[X,Y]/(X^2+Y^2)$, a real domain whose only real point is the origin (Artin's example after SGA~4 X~2.3). The margin also asks: \q{cette restriction est-elle essentielle ???}}

\known The inequality is SGA~4 X~2.4 when $r$ is invertible in $A$, with the cruder exception \q{some residue field orderable} \cite{sga4}. When $r$ is the residue characteristic and $\operatorname{char} K = 0$, it follows for $A$ excellent henselian from Gabber--Orgogozo~\S3 \cite{gabberorgogozo}. We found nothing on the finer exception~(b), on the regular case, on the catenary refinement, or on $r$ the residue characteristic for $A$ neither excellent nor henselian. Serre's \emph{Cohomologie galoisienne} II~\S4 and Scheiderer's \emph{Real and Étale Cohomology} \cite{scheiderer}, where cohomological dimension in the presence of orderings is treated, remain to be read.

\ideas Artin's proof of X~2.4 goes by induction on $\dim A$, down to a henselian local ring of dimension one through a height-one prime, where X~2.3 and 2.2 apply \cite{sga4}. The finer exception~(b) says that what matters is the virtual cohomological dimension of the real residue fields, which suggests running the same induction with $\cd_2$ replaced by $\cd_2$ of $k(x)(\sqrt{-1})$ at the real points. The page's exception list covers $r$ equal to the residue characteristic in mixed characteristic, which SGA~4 X~2.4 does not.

\leanst Not attempted (Galois cohomology of fields is in Mathlib only in low degree).

\settle A reference for (b) and the two refinements, or an example of a regular local domain where the inequality fails at $r = 2$.

\begin{problem}[Hodge types of the deleted neighbourhood of a point]\label{p:box}
Let $Y$ be a complex algebraic variety of dimension $n$, $j : X \hookrightarrow Y$ an open immersion with $X$ smooth and dense, $y \in Y \smallsetminus X$, and $H^i = H^i(X \cap B_y) = (R^i j_*\mathbb Q)_y$ with its mixed Hodge structure (Durfee \cite{durfee} when $(Y \smallsetminus X) \cap B_y = \{y\}$; in general through Saito's mixed Hodge modules \cite{saito}). Then
\[
h^{p,q}(H^i) \ne 0 \;\Longrightarrow\; p, q \in [\max(0, i-n+1),\ \min(i, n)],
\]
so the weights lie in $[0, 2i]$ for $i \le n-1$ and in $[2(i-n+1), 2n]$ for $n \le i \le 2n-1$ (and $H^i = 0$ for $i \ge 2n$), and the level $\max|p-q|$ is at most $\min(i, 2n-1-i) \le \min(i, n-1)$. Is this in the literature, and what is the right relative statement for a proper morphism?
\end{problem}

\src{\fac{12}, pp.~115--119, \emph{Généralités (sous-groupes de Galois motiviques)}, inventory dating 1964--[c.~1971]. Page~119 boxes \q{poids $\le 2\inf(i,n)$, niveau $\le \inf(i, n-1)$} for $X$ regular, $F = \mathbb Q_\ell$ and $f$ an immersion with $n = \dim \overline{f(X)}$, in his yoga of weights and levels of $\ell$-adic sheaves. The same page asks, with a question mark, for a relative version with $n = d_y(f)$ and level $\le \min(i,n)$. The typescript says of its own conclusions (p.~96): \q{la portée des \q{résultats} de notre travail est nécessairement sujette à caution, puisqu'ils dépendent de conjectures non prouvées à l'heure actuelle} (the Weil conjectures and resolution of singularities).}

\ctx The page's argument is a Leray spectral sequence for a factorisation of $f$ as an open immersion with normal crossings complement followed by a proper map, conditional on resolution, purity and the Weil conjectures, with tables of weights and levels. The Hodge-theoretic statement above is our translation.

\known Deligne bounds the Hodge types of the cohomology of a variety \cite[8.2.4]{hodge3}; Durfee constructs the mixed Hodge structure on $H^i(X \cap B_y)$ and its weight spectral sequence; Durfee--Saito, Steenbrink, Navarro Aznar and El~Zein bound weights; for isolated singularities the box is implicit in the literature on links: it follows from semipurity and duality as in \cite[Th.~2.6, Def.~2.7]{friedmanlaza}. We did not find the box, or the level bound, for arbitrary $y$.

\ideas Our argument: take a resolution $\pi : \tilde Y \to Y$, an isomorphism over $X$, with $D = \pi^{-1}(Y \smallsetminus X)$ simple normal crossings; put $F = \pi^{-1}(y)$ and $F_q = F \cap D^{(q)}$. Then $\dim F_0 \le n-1$ and $\dim F_q \le n-q$. The weight spectral sequence $E_1 = H^{k-q}(F_q)(-q) \Rightarrow H^k(X \cap B_y)$ (Durfee's, in the case he treats; we have not checked its form in the general case) and Hodge~III 8.2.4~(i), (ii) for the compact $F_q$ put every $E_1$ term in the box. Locally the lower end moves up by one because $\pi^{-1}(y)$ has dimension at most $n-1$. The cone over a smooth projective $V$ of dimension $n-1$ with $h^{i,0}(V) \ne 0$ attains level $i$ for $i \le n-1$ and $n-1$ at $i = n$. For a proper $f$ with fibre dimension $d$ the gain fails, since a smooth family of curves already has $H^1$ of level $1 = d$; so the relative question of page~119 cannot keep $n-1$ if $d_y(f)$ is the fibre dimension.

\leanst Not attempted: no Hodge structures on cohomology in Mathlib.

\settle A reference for the box; a relative statement interpolating between the immersion case and the smooth proper case; its $\ell$-adic form with level read as divisibility of Frobenius eigenvalues.

\begin{problem}[An $\Ext^2$ that vanishes when $2$ is invertible: beyond smooth coefficients]\label{p:ext2}
On the fppf site of a base on which $2$ is invertible, is $\mathcal Ext^2_{\mathbb Z}(G, M) = 0$ for every finite locally free commutative group scheme $G$ and every quasi-coherent $M$, not only for the smooth coefficients covered by Breen's theorem below?
\end{problem}

\src{\fac{129}, p.~95, \emph{Notes séminaire Montréal}, inventory dating [1970]. A parenthetical NB after a theorem on Ext over truncated rings: \q{$\Ext^2_{\mathbb Z}(G, M) = 0$ si 2 inv., pour tt $G$ fini loc.\ libre}, with the remark that this simplifies the deformation theory of Barsotti--Tate groups for $p \ne 2$. The word \q{inv.} is uncertain in the transcription, and the page gives no proof.}

\known Breen proved, five years later, that $\mathcal Ext^i(G, H) = 0$ for $1 < i < 2p-2$, for $G$ finite flat and $H$ smooth commutative quasi-projective, when every residue characteristic is at least $p$ and multiplication by every prime $q < p$ is an automorphism of $G$; and that $\mathcal Ext^{2p-2}(\alpha_p, H) \ne 0$ in characteristic~$p$ for $H = \mathbb G_a, \mathbb G_m$ \cite{breen1975}. With $p = 3$ this gives the NB for smooth coefficients, in particular for vector bundles, when $2$ is invertible; with $p = 2$ it shows $\mathcal Ext^2(\alpha_2, \mathbb G_a) \ne 0$, which supports the reading \q{si 2 inv.}. Illusie places the obstruction to deforming a truncated Barsotti--Tate group over a perfect field in $\Ext^2_{\mathbb Z}(G_0, t_{G_0} \otimes J)$ and says that proving directly that it vanishes is out of reach \cite[p.~2]{illusie2023}; Breen's theorem, which he does not cite, bears on that group.

\ideas Breen's theorem is stated for the étale sheafification of the Ext sheaves. The remaining cases are a quasi-coherent $M$ that is not locally free, and the passage from Breen's sheaves to the global groups on the fppf site, through the local-to-global spectral sequence.

\leanst Not attempted: no fppf cohomology in Mathlib.

\settle Breen's vanishing for an arbitrary quasi-coherent $M$, or a counterexample; and whether it gives Illusie's obstruction directly.

\section{Algebraic cycles}

\begin{problem}[Splitting the critical Lefschetz condition]\label{p:lef}
Let $X$ be smooth projective of dimension $n$ over an algebraically closed field, with a Weil cohomology $H$, a smooth hyperplane section $\varphi : Y \hookrightarrow X$, and $L$ the Lefschetz operator. Write $H^{2i}$ for the cohomology in which codimension-$i$ cycles live, and consider:
\begin{itemize}
\item $A'(X)$: for $2i \le n-1$, a class $x \in H^{2i}(X)$ whose restriction $\varphi^*x$ is algebraic is algebraic;
\item $A''(X)$: for $2i \ge n-1$, every algebraic class of $H^{2i+2}(X)$ is $\varphi_*y$ for an algebraic $y \in H^{2i}(Y)$;
\end{itemize}
and their critical versions $A'^{0}$, $A''^{0}$, with $2i = n-2$ for $n$ even and $2i = n-1$ for $n$ odd. Let $A^0(X)$ be the critical Lefschetz condition: $L^2 : H^{n-2}(X) \to H^{n+2}(X)$ for $n$ even, $L : H^{n-1}(X) \to H^{n+1}(X)$ for $n$ odd, is a bijection on algebraic classes. Under which hypotheses do $A'^{0}(X)$ and $A''^{0}(X)$ imply $A^0(X)$, and is the resulting reduction of the Lefschetz standard conjecture,
\[
C(X) \iff C(Y) + A'^{0}(X \times X) + A''^{0}(X \times X),
\]
correct as stated?
\end{problem}

\src{\fac{16}, p.~49, \emph{Formalisme algébrique des correspondances et algèbres de Poincaré-Hodge}, inventory dating 1966--1967, in a letter of January 1967. The letters use $A$, $C$, $D$ in senses that do not match the later published labels; the conditions above are the reading's restatement. The letter to Coates places the question in the motivic picture (p.~48): \q{The complete yoga then being that we have the category of motives which is semi-simple !}.}

\known The equivalence $C(X) \iff C(Y) + A^0(X \times X)$ and its descent along a chain of hyperplane sections are Kleiman's \cite[Th.~2.13]{kleiman1968}. The splitting of $A^0$ into the two one-sided conditions was not found in Kleiman 1968 and 1994, Grothendieck 1969, Charles 2013, Ramón Marí, or Milne; André's book \cite[ch.~5, \S5.2]{andre} remains to be read.

\ideas The page factors the critical operator through the section. For $n = 2m$: $L^2 = \varphi_* \circ L_Y \circ \varphi^* : H^{2m-2}(X) \to H^{2m-2}(Y) \to H^{2m}(Y) \to H^{2m+2}(X)$; for $n = 2m+1$: $L = \varphi_*\varphi^* : H^{2m}(X) \to H^{2m+2}(X)$. With it (our reconstruction, $n = 2m$): an algebraic $z \in H^{2m+2}(X)$ is $\varphi_*y$ with $y$ algebraic by $A''^{0}$; by weak and hard Lefschetz, $y = L_Y\varphi^*x$ for a unique $x$, and $z = L^2x$; if $L_Y^{-1}$ preserves algebraic classes on $Y$, $\varphi^*x$ is algebraic, and then $x$ is algebraic by $A'^{0}$. So the splitting holds given the Lefschetz condition on the section in the relevant degree. In the displayed equivalence it is applied to $T = X \times X$, whose hyperplane section is not $Y \times Y$; whether the hypothesis on the section is supplied by the induction, or must be added, is the point to check.

\leanst Not attempted: no Weil cohomology in Mathlib.

\settle A reference for the splitting, or a precise statement of the hypotheses under which it holds, applied to $X \times X$.

\section{Abelian varieties and reductive groups}

\begin{problem}[Polarised abelian surfaces with even Euler characteristic]\label{p:surf}
Let $A/S$ be an abelian scheme of relative dimension~$2$ with dual $B$, and $\delta$ a section of $\mathrm{NS}_{A/S}$ with $\chi(\delta)$ even. Is the torsor $\Pic^\delta_{A/S}$ under $B$ trivial? Over a field $k$, this says that the Poonen--Stoll class $c_\delta \in H^1(k, B)$ vanishes \cite{poonenstoll}.
\end{problem}

\src{\fac{59}, pp.~11--13, \emph{Construction des faisceaux inversibles sur schémas de Picard}, inventory dating [années 1960--1970]. Page~11 gives a canonical isomorphism $\Pic^{\delta'}_{B/S} \simeq \tilde\delta'_*\Pic^\delta_{A/S}$; page~13 composes it with its analogue and concludes, for $n = 2$ and $\chi(\delta)$ even, that $\Pic^\delta$ is trivial. That last line rests on an uncertain \q{si} and an illegible word.}

\known Poonen and Stoll define $c_\lambda$ and prove $2c_\lambda = 0$, with no criterion in terms of $\chi(\lambda)$ or the dimension \cite[\S4]{poonenstoll}. We found no such criterion elsewhere.

\ideas The reading thought the conclusion depends on compatibilities of trivialisations that the folder asks for and never proves. Our argument, unchecked, needs none for mere triviality. The map $L \mapsto \det(\hat L)^{\pm1}$, with $\hat L$ the Fourier--Mukai transform, is a morphism of fppf sheaves $\Pic^\delta_{A/S} \to \Pic^{\varphi(\delta)}_{B/S}$, equivariant along $\varphi(\delta)^\sim$ by the theorem of the square. The same construction on $B$ goes back to $\Pic^{\varphi\varphi(\delta)}_{A/S}$, and in dimension~$2$, $\varphi\varphi(\delta) = \delta$. The composite $B \to B$ is multiplication by $\mp\chi(\delta)$. Hence $(\chi(\delta) \pm 1)\,[\Pic^\delta] = 0$, and $2\,[\Pic^\delta] = 0$ through $L \mapsto L \otimes [-1]^*L$. Both $\chi \pm 1$ are odd when $\chi$ is even, so the class vanishes whatever the signs. Whether the trivialisation obtained is canonical is a separate question.

\leanst Not attempted: no abelian schemes in Mathlib.

\settle A specialist's check of the argument. One abelian surface over a number field with a polarisation of type $(1,2)$ and $c_\lambda \ne 0$ refutes the statement. A sharp consequence to test (our heuristic, unchecked): for a genus-$2$ curve $C/k$ whose Jacobian has real multiplication by $\mathcal O_D$, $D \equiv 1 \bmod 8$, defined over~$k$, the theta class is a sum of two classes of even $\chi$, so the statement would force $\Pic^1_C$ to have a $k$-point.

\begin{problem}[Kostant's tori in apposition: are the Borels opposite?]\label{p:kostant}
Let $G$ be a complex simple adjoint group with Coxeter number $h$ and Weyl group $W$, and let $T$, $T'$ be maximal tori in apposition in Kostant's sense \cite{kostant}, with Borel subgroups $B \supset T$ and $B' \supset T'$. Are $B$ and $B'$ always in general position, that is, is $B \cap B'$ a maximal torus? If not for all such quadruples $(T, B, T', B')$, is there a distinguished conjugacy class of quadruples, stable under the outer automorphisms?
\end{problem}

\src{\fac{123}, pp.~21--22, \emph{[Autour de Kostant]}, letter to Kostant of 22 October 1969: \q{is it true that for some, or for all, such quadruples, $B$ and $B'$ are \q{in general position} [\dots] ? Is there in any sense a \emph{distinguished} conjugacy class of such quadruples [\dots] ?} The letter states without proof that the quadruples form $|W|^2/(h\,z\,\varphi(h))$ conjugacy classes, $z$ the order of the centre of the simply connected cover, and (in a footnote in his hand) that $N(T) \cap N(T')$ is an extension of $(\mathbb Z/h)^*$ by $\mathbb Z/h \times Z(\tilde G)$.}

\known The count, which this note listed as open until 26 October 2026, follows from standard theory (our argument, unchecked by a specialist). Let $x = \rho^\vee(e^{2\pi i/h})$, the principal element of $T$, and $T'$ the centraliser of a cyclic element $E' = \sum_i c_i e_{\alpha_i} + c_0 e_{-\theta}$. An element $t \in T$ normalising $T'$ sends $E'$ into $\mathfrak h' \cap \mathfrak g^{\mathrm{Ad}(x) = \zeta}$, which is the line $\mathbb C E'$, so $\alpha_i(t) = \theta(t)^{-1} = \lambda$ with $\lambda^h = 1$: $T \cap N(T') = \langle x \rangle$, of order $h$. The image of $N(T) \cap N(T')$ in $W$ is the set of $w$ with $w(x) \in \langle x \rangle$: the stabiliser of $x$ is the group $\Omega \cong P^\vee / Q^\vee$ of order $z$ (in the adjoint group $x$ has a disconnected centraliser), $x$ is conjugate to $x^k$ exactly for $k$ prime to $h$ \cite{kostant,springer}, and each such $w$ lifts into $N(T')$, because regular semisimple elements of the eigenspaces $\mathfrak g^{\zeta^j}$, $j$ prime to $h$, have the same centralisers as those of $\mathfrak g^{\zeta}$ and their Cartan subspaces are conjugate under $Z_G(x)^\circ = T$. Hence $|N(T) \cap N(T')| = h\,z\,\varphi(h)$; the stabiliser of a quadruple is $T \cap T' = 1$, and the count follows. A computer count over $\mathbf F_q$ (\nolinkurl{computations/kostant}) agrees: 2, 4 and 12 orbits for $A_2$, $B_2$ and $G_2$. For general position we found nothing.

\ideas For $A_2$ all $36$ pairs $(B, B')$ are opposite: the eigenvectors of $E'$ are the Fourier vectors $(1, \omega^k, \omega^{2k})$, and the standard flag is transverse to every flag they span. In general the question is the Bruhat cell of $(B, B')$ for $B'$ running over the $|W|$ Borels through $T'$: whether it is always the open cell.

\leanst Not attempted in Lean. The computer count above can be extended to the relative position of $B$ and $B'$ in $B_2$ and $G_2$, which we have not done.

\settle The relative position of $B$ and $B'$ for all quadruples, by an argument or by computation in each type.

\section{Topoi and stacks}

\begin{problem}[Two classifying topoi that are not coherent]\label{p:noncoh}
The theories of commutative rings with irreducible spectrum (nilradical prime) and of rings with zero-dimensional spectrum are geometric. The page says their classifying topoi, subtopoi of the coherent classifying topos of rings, are not coherent, their presenting topologies on affine schemes of finite type over $\mathbb Z$ not being quasi-compact. Is that right, and do these topoi have enough points?
\end{problem}

\src{\fac{161-2}, pp.~39, 41--42, \emph{Algèbre universelle}, inventory dating [c.~1963--1973]. Page~39: \q{Cette topologie n'est pas quasi-compacte}; page~42, margin: \q{topologie non cohérente ! A-t-elle assez de pts ??}; page~41, margin: \q{encore un topos non cohérent (bien que sous-topos d'un topos cohérent)}.}

\known Locally coherent toposes have enough points (Deligne, \cite[VI, 9.0]{sga4}); these theories are geometric but, as presented, use the infinite disjunction hidden in \q{nilpotent}. We found neither non-coherence nor the existence of points stated for them, in Johnstone \cite{elephant} or in Caramello's criteria for De Morgan classifying toposes, which take quotients of the theory of rings as examples \cite[\S5]{caramello}.

\ideas Non-coherence needs more than the presentation: one must show that no coherent axiomatisation exists, for instance by exhibiting a cover in the classifying topos with no finite subcover. For points, Barr's theorem gives a Boolean cover, not points; a model-theoretic argument, or an explicit set of local rings detecting the theory, is needed.

\leanst Not attempted: classifying topoi are not in Mathlib.

\settle A proof or reference for non-coherence, and an answer to the question about points.

\begin{problem}[Foliations and their quotient topoi]\label{p:fol}
Is there an equivalence between the category of pairs $(X, R)$, with $X$ a topological space and $R$ a local equivalence relation on $X$, locally open with locally connected fibres, and the category of triples $(X, \mathcal Y, f)$, with $f : \Top(X) \to \mathcal Y$ a locally open geometric morphism with locally connected fibres to a \q{multiplicité topologique}, a $2$-category whose Hom-categories are discrete?
\end{problem}

\src{\fac{119}, pp.~144--145, dated \q{Mars 1986} on p.~144. Under \q{Th}: \q{On a $C \rightleftarrows C'$}, both arrows marked $\approx$. No proof. A third condition on $f$ is partly illegible; the reading suggests it may be a simple connectedness of the fibres.}

\known SGA~4 IV, Exercise~9.8.3 j)--k), gives the functor $(X, R) \mapsto (\Top(X) \to \Top(X/R))$ \cite{sga4}; Rosenthal states that the topos $\mathrm{sh}(X; r)$ exists for strictly open $r$ (Th.~3.9, proof in his thesis) and proves it is an étendue for strictly regular $r$ \cite{rosenthal}; Kock and Moerdijk show that every étendue is of the form $\mathrm{sh}(M, r)$, construct the inverse $a \mapsto \mathrm{Lker}(a)$ and show the round trips are identities under simple connectedness or kernel-pair conditions \cite{kockmoerdijk1992}, and study monodromy versus holonomy \cite{kockmoerdijk1996}. We found no statement of an equivalence of categories.

\ideas The two functors are there; the problem is the right conditions on both sides so that they are mutually inverse, and which quotient topos the page means: the topos of $R$-invariant sheaves (sheaves on a monodromy groupoid) or the holonomy quotient, which Kock and Moerdijk distinguish.

\leanst Not attempted: no geometric morphisms of topoi in Mathlib.

\settle A correct statement with proof, or Moerdijk's Zaragoza lectures \cite{moerdijk1993}, which we could not read and where it might be.

\begin{problem}[Realising the symmetry invariant of a Picard stack]\label{p:pic}
Let $M$, $N$ be abelian sheaves on a topos, and ${}_2N$ the $2$-torsion of $N$. Is every homomorphism $M \to {}_2N$ the symmetry invariant of a Picard stack pinned by $M$ and $N$?
\end{problem}

\src{\fac{134-1}, p.~5, letter to Deligne, Lodève, 27 August 1974 (Breen's copy, forwarded by Deligne): \q{Je sais prouver (sauf erreur) que tout hom.\ $M \to {}_2N$ provient d'un champ de Picard convenable (épinglé par $M$, $N$) (a priori l'obstruction est dans $\Ext^3(X; M, N)$, mais un argument \q{universel} prouve qu'elle est nulle)}. No argument is written. On p.~6 the letter asks his own version of this problem: \q{Question : Le triangle exact (T) et la suite exacte (*) sont-ils connus par les compétences (Quillen, Breen, Illusie\dots) ? Connaissent-ils des variantes \q{supérieures} ?} (are they known to the experts, and do they know higher variants?).}

\known Strict Picard stacks are classified by $\Ext^2$ (Deligne, \cite[XVIII, 1.4]{sga4}); Breen classifies Gr-stacks by $H^3$ and gives the universal coefficient spectral sequence $\Ext^p(H_q(B), A) \Rightarrow H^{p+q}(B, A)$ \cite[\S6]{breen1999}. No printed source we found states the realisation, or refutes it on a topos.

\ideas Two sketches disagree. The reading realises the invariant by hand for $M$ free, on the split stack $M \times BN$ with an antisymmetric form $b$ such that $b(e_i, e_i) = q(e_i)$, and passes to general $M$ through a resolution. We doubt that step, which we have not checked: the construction depends on a basis. On the classifying topos of abelian groups, a Picard stack pinned by the generic group $U$ and $U/2U$, with the projection as invariant, would make Moore spectra truncated above $\pi_1$ functorial in finitely presented abelian groups; but the stable Moore-space category on finitely generated groups is a non-split linear extension, split on cyclic groups \cite[Rem.~1.5.5]{baues}. If that translation is right, the obstruction does not vanish universally. The statement does hold over a point and for $M$ constant and free.

\leanst Not attempted: no stacks on a topos in Mathlib.

\settle The differential $\operatorname{Hom}(M/2, N) \to \Ext^3(M, N)$ of the universal coefficient spectral sequence on the classifying topos of abelian groups, or Grothendieck's \q{universal} argument written out.

\begin{problem}[Morphisms of ramification domains are fully faithful]\label{p:domram}
Over a scheme $S$, a \q{domaine de ramification} is a stack of multigaloisian categories over the étale site of $S$ with a continuous geometric realisation into relative schemes, subject to axioms R1 and R2 (pp.~190--193). For a domain $(R, r)$ defined by a scheme $Z$ with a finite group $G$ of operators, is $\operatorname{Hom}((R, r), (R', r'))$ equivalent to the category of pairs $(P', \alpha)$, $P'$ a galoisian object of $R'_S$ of group $G$ and $\alpha$ a $G$-isomorphism $r'(P') \simeq Z$, so that every homomorphism of domains is fully faithful?
\end{problem}

\src{\fac{29}, p.~197, \emph{Groupe fondamental [autour de SGA~4 et SGA~7]}, inventory dating 1959--1969: Proposition~6, stated without proof, with the corollary \q{nécessairement pleinement fidèle}. The axiom R2 of pp.~190--193 is itself not verified by the edition. Earlier in the folder (p.~24) he notes where removing the auxiliary choices leads: \q{il semble qu'on soit infailliblement conduit à prendre comme catégorie [\dots] la catégorie [\dots] elle-même qu'on se propose de construire et d'étudier !}, comparing it with sites and their topoi.}

\known Nothing found. Tonini's work on stacks of ramified Galois covers \cite{tonini2013,tonini2017} is where an equivalent statement would sit; neither surfaced in our search, which is not the same as their not containing it.

\ideas The corollary says that a ramification domain maps into another only as a subdomain. The proof should reduce to the case $(R, r)$ given by $(Z, G)$, where a morphism is determined by the image of the universal $G$-torsor.

\leanst Not attempted: no étale site or stacks of categories in Mathlib.

\settle A proof of Proposition~6 from R1--R2, or a reference.

\section{Moduli of curves and Galois actions}

\begin{problem}[The Galois image, seen on one curve]\label{p:gal}
Let $X$ be an anabelian curve over a field $K$ of finite type over $\mathbb Q$, and $\hat\Pi_1(X, I)$ its profinite fundamental groupoid on a finite set $I$ of base points. The page lists necessary conditions for a vertex-fixing automorphism $\Theta$ of $\hat\Pi_1(X, I)$ to come from $\Gal(\bar K/K)$: it lifts to $\hat\Pi_1(X', I')$ for every finite étale cover $X'$ of $X$ defined over $K$; it preserves the kernels of $\hat\pi_1(X', x') \to \hat\pi_1(Y, f(x'))$ for every $K$-morphism $f : X' \to Y$; it is compatible with morphisms between covers; and it satisfies the loop condition at the cusps with a single multiplier. Are these conditions sufficient?
\end{problem}

\src{\fac{142}, pp.~94--95, 98--101, \emph{Action de Galois sur Teichmüller}, inventory dating [from 1978]. Page~95 conjectures the conditions sufficient \q{lorsque $X$ est (disons) une courbe algébrique anabélienne, avec $K$ extension de type fini de $\mathbb Q$}; page~99 calls the loop condition \q{sûrement la plus importante de toutes, et de très loin !}. Pages~98--99 are largely illegible.}

\known Ihara's question and the Oda--Matsumoto conjecture characterise the Galois group as the automorphisms of the geometric fundamental group functor on all varieties compatible with all morphisms; Pop proved it in that form, and finite tripod variants, built from varieties of dimension greater than one with finitely many dominant rational maps to $\mathbb P^1 \smallsetminus \{0,1,\infty\}$ \cite[Def.~2.2, Th.~2.9]{pop2019}. Hoshi--Mochizuki and others give injectivity for a single hyperbolic curve; Collas proposes a combinatorial model $GT^{NF}$ as a candidate for $G_{\mathbb Q}$, left to forthcoming work \cite[\S4.3.3]{collas}. We found no characterisation attached to one curve and its covers, with base points and the inertia condition.

\ideas Pop's variants need a base of dimension greater than one, so they do not cover the covers of one curve; that is what the page's loop condition at the cusps would have to supply. Here \q{anabelian curve} is the page's word for a hyperbolic curve.

\leanst Not attempted: no étale fundamental group in Mathlib.

\settle A proof of sufficiency, a counterexample, or a published variant that implies it.

\begin{problem}[The deleted neighbourhood of a normal crossings divisor]\label{p:nc}
Let $X$ be a smooth complex analytic space (or orbifold) with a normal crossings divisor $\Theta$, and $V^*_{\Theta,X}$ the deleted tubular neighbourhood of $\Theta$. Is $V^*_{\Theta,X}$ the homotopy colimit of the semi-simplicial system of deleted tubular neighbourhoods of the strata of the unfolding of $\Theta$, so that its fundamental groupoid is the amalgamated sum $\tilde\Pi_1 \leftarrow \tilde\Pi_{2,1} \rightarrow \tilde\Pi_2$ of the groupoids of the torus bundles over the open strata of codimension $1$ and $2$? And hence the same for $\Pi_1(X \smallsetminus \Theta)$ whenever $X \smallsetminus \Theta$ and its end have equivalent fundamental groupoids?
\end{problem}

\src{\fac{147}, pp.~148--151, \emph{Structure à l'infini des $M_{g,\nu}$}. The folder is undated; the inventory dates its group (141--149) [from 1978]--1983. We date it to the winter of 1982--1983, because its formalism is the one Grothendieck describes to Illusie in a letter of 22--25 January 1983 (IHÉS, fonds Grothendieck--Illusie, 1ILL35, p.~27) and in a \emph{rapport d'activité} of which he sent Illusie a copy on 29 March 1983 (1ILL41, p.~4; copy in \fac{119}, pp.~141--142). Page~149 states a \q{Th.\ de recollement} as an equivalence of topoi and corrects it at once: \q{en fait, c'est une équivalence d'homotopie}. Pages~150--151 give the corollaries, with a purity argument for neglecting the strata of codimension~$3$. No proof is written.}

\ctx The folder describes the boundary of $\overline M_{g,\nu}$ by stable graphs, lists the strata of modular dimension at most~$2$, and writes on its p.~29 (archive p.~59) that a more or less complete understanding of the discrete $\pi_1$ in these small cases \q{donnera la clef pour une compréhension complète dans le cas général}: the two-level principle of the Teichmüller tower, about a year before the \emph{Esquisse d'un programme}. Applied to $X \smallsetminus \Theta = M_{g,\nu}$, the corollaries would express the mapping class group through the strata of codimension $1$ and~$2$; that application is the reading's, not the page's.

\known The two-level principle for the discrete tower was proved by Hatcher, Lochak and Schneps \cite{hls}. For the general gluing statement, the places to look are the real oriented blow-up and the Kato--Nakayama space of the log structure, whose boundary is a torus-bundle model of the deleted neighbourhood, and Looijenga's and Boggi's work on the boundary of $\overline M_{g,n}$. We did not find the statement in this form.

\ideas The purity step reads: removing a subset of real codimension at least~$3$ does not change $\pi_1$. For $M_{g,\nu}$, the hypothesis of the last corollary (end and interior with the same $\Pi_1$) should follow from the simple connectivity of the curve complex (Harer).

\leanst Not attempted: no homotopy colimits of spaces or stratified spaces in Mathlib.

\settle A reference for the gluing theorem, or a proof, and the check of the hypothesis for $M_{g,\nu}$.

\section{Pseudoline arrangements}\label{sec:pseudo}

A pseudoline arrangement in the real projective plane is \emph{bigonal along} one of its pseudolines $D$ when $D$ carries exactly two vertices of the arrangement, so that every other pseudoline meets $D$ at one of them. In the page's notation, $\mathrm{II}_4$ is four pseudolines in general position, $\mathrm{II}_5$ five in general position, $\mathrm{II}'_{3,2}$ a pencil of three and a pencil of two without the line joining their centres, and $\mathrm{III}_6$ the complete quadrangle (the six lines joining four points, no three collinear).

An arrangement is \emph{bigon-free} when none of its pseudolines is bigonal, and four pseudolines are in \emph{general position} when no three of them pass through one point. The page's boxed dichotomy, which this note listed as an open problem until 26 October 2026, has a short proof, and so does the numbering lemma of page~28; we record it because the problem below depends on it.

\begin{proposition}[Grothendieck, \fac{80}, pp.~20, 24, 28; the proof is ours]
A bigon-free arrangement $\mathcal A$ either is $\mathrm{II}_4$ or contains $\mathrm{II}_5$, $\mathrm{II}'_{3,2}$ or $\mathrm{III}_6$; and its pseudolines can be numbered, starting from a $\mathrm{II}_4$ or a $\mathrm{III}_6$, so that every partial arrangement is bigon-free.
\end{proposition}

\begin{proof}
\emph{Four pseudolines in general position.} $\mathcal A$ is not a near-pencil, since each pencil line of a near-pencil carries only the centre and its point on the transversal; so it has a simple crossing $p = a \cap b$ \cite[Th.~3.6]{kellyrottenberg}, see also \cite[Th.~1.3]{bokowskipokora}. Since $a$ is not bigonal it carries vertices $q_1, \dots, q_k$ besides $p$, with $k \ge 2$. Take $c \ne a$ through $q_i$ and $d \ne a$ through $q_j$, $i \ne j$. Then $\{a, b, c, d\}$ is in general position as soon as $c \cap b \ne d \cap b$: indeed $c, d$ miss $p$, meet $a$ at different points, and $c \cap d$ lies neither on $a$ nor on $b$. If this fails for every choice, all these pseudolines meet $b$ at one point $r$; two pseudolines meet only once, so each $q_i$ lies on $a$ and on a single other pseudoline $c_i$, all through $r$, and every pseudoline meets $a$ at $p$ or at some $q_i$. Then $\mathcal A = \{a, b, c_1, \dots, c_k\}$ and $b$ carries only $p$ and $r$, so it is bigonal, which is excluded.

\emph{The dichotomy.} Let $Q \subset \mathcal A$ be four pseudolines in general position, and $\ell \in \mathcal A \smallsetminus Q$. If $\ell$ passes through none of the six vertices of $Q$, then $Q \cup \{\ell\}$ is a $\mathrm{II}_5$; if through exactly one, a $\mathrm{II}'_{3,2}$. Otherwise $\ell$ passes through exactly two, which must be an opposite pair $\{x \cap y, z \cap w\}$ of the quadrilateral $Q = \{x,y,z,w\}$ (two vertices on a common side would force $\ell$ to be that side). Call such an $\ell$ a diagonal of $Q$; there are at most three. If $\mathcal A$ contains neither $\mathrm{II}_5$ nor $\mathrm{II}'_{3,2}$, then $\mathcal A$ is $Q$ plus some diagonals. One diagonal alone would be bigonal; two diagonals together with $Q$ form a $\mathrm{III}_6$ (the four triple points are the vertices of the quadrangle and the six pseudolines its sides); with three, $\mathcal A$ contains a $\mathrm{III}_6$ again.

\emph{The numbering.} Adding pseudolines never removes a vertex from a pseudoline already there, so it suffices that each pseudoline, when added, meets the earlier ones in at least three points. A pseudoline meets the six sides of a $\mathrm{III}_6$ in at least three points: its vertices have multiplicity at most three, and two triple points share a side, so two vertices cover at most five sides. So if $\mathcal A \supseteq \mathrm{III}_6$, start there and add the rest in any order. Otherwise $\mathcal A$ contains at most one diagonal $\delta$ of $Q$: start from $Q$, add every pseudoline other than $\delta$ (each meets $Q$ in at least three points), and add $\delta$ last, when it has all its vertices in $\mathcal A$, at least three.
\end{proof}

This agrees with a computer enumeration of all arrangements of at most nine pseudolines (\nolinkurl{computations/bigon}). The dichotomy is also checked in Lean for combinatorial maps (generalised maps of Euler characteristic~$1$ whose flags are labelled by the pseudolines; \nolinkurl{lean/Grothendieck/Folder80Maps2.lean}, \texttt{map\_dichotomy}); the numbering and the passage from a topological arrangement to its map are not formalised. We have not found the statement in print; it is elementary, and probably folklore among people who work on arrangements.

\begin{problem}[A bigon-free arrangement is determined by its skeleton]\label{p:skel}
Is a combinatorial arrangement of pseudolines in the real projective plane without bigonal pseudolines determined by its $1$-skeleton together with the decomposition of the edges into pseudolines, functorially: is the functor from bigon-free arrangements to their decomposed skeleta fully faithful, so that every automorphism of the decomposed skeleton comes from the arrangement?
\end{problem}

\src{\fac{80}, pp.~39--54. The statement of page~54 (\q{En résumé}) is clean. It composes a full faithfulness theorem (p.~39) for arrangements with an orientation of every non-concurrent triple (pp.~41--49, conditions (A1)--(A4)) with a canonical reconstruction of these orientations from the skeleton (pp.~49--54, (A5), (A6), with the complete quadrangle as residual case). Conditions (A1)--(A6) are said to hold \q{bien entendu} and not proved.}

\known Line arrangements in the projective plane are determined up to isomorphism by their circular sequences up to local equivalence (Goodman--Pollack), and the graph of a simple projective arrangement of at least four pseudolines is $4$-connected \cite[Th.~5.2.2, 5.5.1]{felsnergoodman}. Embeddings of graphs in the projective plane are unique at face-width at least~$3$ and are classified at face-width $2$ and $3$ \cite{vitray,mohar1995,mohar-robertson-vitray}. For simple arrangements the question reduces to the plane, by our argument, unchecked by a specialist: deleting one labelled pseudoline leaves a simple Euclidean arrangement; its graph, with one vertex added and joined to the vertices of degree less than four, is $3$-connected and planar \cite{boseeverettwismath}, so by Whitney's theorem every isomorphism of such graphs (which preserves degrees) comes from a homeomorphism of the sphere, and the pseudolines are the paths through opposite edges \cite[Lemma~1]{eppstein}. This simple case is also the one proved in Lean (see below). We found no statement for arrangements with multiple points, which is the open part.

\ideas Our pointer, unchecked: if the graph is $3$-connected and its embedding has face-width at least~$3$, these uniqueness theorems give the statement, without even the decomposition into pseudolines, since an automorphism then maps facial walks to facial walks. A non-contractible curve meets every pseudoline, so face-width~$2$ means two points covering all of them; for a bigon-free arrangement this happens exactly for two pencils with no pseudoline through both centres, such as $\mathrm{II}_4$ and $\mathrm{II}'_{3,2}$, which then need a separate argument, as does $3$-connectivity in the non-simple case.

\leanst Proved for simple arrangements, in the combinatorial maps of page~60 with its axiom that each pseudoline goes straight through each vertex (\nolinkurl{lean/Grothendieck/Folder80Maps3.lean}). The functor is faithful (\texttt{faithful}), and a skeleton isomorphism lifts exactly when it preserves the face paths of length three (\texttt{lift\_iff}, page~41). The orientation of a triangle is read on the skeleton when a fourth pseudoline avoids its vertices ((A5), \texttt{triangle\_iff}, \texttt{SkelIso.triangle\_map}). At vertices of any multiplicity, two edges leaving the ends of an edge on the same side cut out a triangle that bounds, and on opposite sides one that does not ((A4), \texttt{side\_bounds}, \texttt{side\_not}). When every vertex lies on exactly two pseudolines, every isomorphism of decomposed skeleta lifts (\texttt{full\_simple}). Not formalised: the cyclic order at a vertex of three or more pseudolines ((A1)--(A3)), (A6) and the complete quadrangle.

\settle The case with multiple points: a proof of (A1)--(A3) at a vertex of three or more pseudolines, of (A6) and of the complete quadrangle case, or the face-width bound with a uniqueness theorem.

\section{The large conjectures, and how the notes approach them}\label{sec:big}

The fourteen problems above may not be in the literature. Seven conjectures of the notes are, under Grothendieck's name, and the companion paper lists where each first appears in the fonds \cite{huacicm}. This section reads the same pages for how he meant to prove them. For each conjecture it gives the statement in today's form, the pages, a draft of the route the pages take, where the route stops, and what is known now. The drafts are the edition's, made from the transcriptions with a language model. They report what the pages do, and say so when a step comes from the modernised reading instead. None of the seven is solved here, and no route is complete on the page.

\begin{bigconj}[The standard conjectures]\label{c:std}
Let $X$ be smooth projective of dimension $n$ over an algebraically closed field, with a Weil cohomology $H$ and the Lefschetz operator $L$ of an ample class. \emph{Lefschetz type}: the inverse of $L^{n-i} : H^i(X) \to H^{2n-i}(X)$ is induced by an algebraic correspondence. \emph{Hodge type}: on the primitive algebraic classes of $H^{2i}(X)$, $2i \le n$, the form $(-1)^i \langle L^{n-2i}x, x\rangle$ is positive definite.
\end{bigconj}

\src{\fac{14}, pp.~43--46, letter to Serre of 27 August 1965, printed in their correspondence; \fac{14}, pp.~105--113; \fac{16}, pp.~32--54, with letters to Coates (pp.~43--49), inventory dating 1966--1967; \fac{12}, pp.~158--166; \fac{15}, p.~169.}

\entry{The route.} The 1965 letter fixes the order of the argument. A Lefschetz condition on cycles gives a Künneth formula for cycles, hence characteristic polynomials of algebraic endomorphisms with integer coefficients, hence the Weil conjectures except the absolute values of the eigenvalues; the positivity gives the absolute values, and the two together make the motives semisimple. He says he can derive the full Lefschetz condition from a weak, restriction-type one, provided it holds over base fields that are not algebraically closed. Folder~16 then sets up an induction along hyperplane sections $\varphi : Y \hookrightarrow X$. The identity $\Lambda_Y = \varphi^*\Lambda_X^2\varphi_*$ (p.~39) gives $C(X) \Rightarrow C(Y)$, with $C$ the Lefschetz-type conjecture; conversely, \q{Moyennant $C(Y)$, $C(X)$ équivaut à la condition} that $\Lambda_X\varphi_*$ be algebraic (p.~41), which is the Lefschetz condition on $X \times Y$ in the critical degree only. Iterating down a chain of sections reduces $C$ for all varieties to that critical condition for all varieties (pp.~44--45); Problem~\ref{p:lef} is the next step, which splits the critical condition into two one-sided ones (p.~49). A letter to Coates gives a second reduction (pp.~47--48): $C(X)$ follows from integral characteristic polynomials and the existence of \emph{some} algebraic isomorphism $H^{2n-i} \to H^i$, by Cayley--Hamilton. For positivity the route goes through algebras: the rings $\operatorname{End}(h^i(X))$ should carry a positive involution, hence be semisimple, hence the motives form a semisimple abelian category (\fac{12}, p.~158; \fac{16}, p.~48); and one line proposes \q{Positivité à la Hodge en se ramenant à des questions de diviseurs} by induction (\fac{15}, p.~169).

\entry{Where it stops.} The letter of 1965 misses a way of moving a cycle to infinity in the affine complement $X \smallsetminus Y$, which the weak-to-full step needs. The letter to Coates says of the induction: \q{Thus the main problems arise already for 1-cycles on threefolds} (\fac{16}, p.~46). The 1965 conditions are stated modulo algebraic equivalence, and in that form they are false (Griffiths), as the modernised reading notes. For positivity there is only the one line of folder~15.

\entry{Today.} Open in general. The Lefschetz type is known for curves, surfaces and abelian varieties \cite{kleiman1968}, and the Hodge type in characteristic~$0$ by the Hodge--Riemann relations. Deligne proved the Weil conjectures without them.

\begin{bigconj}[The generalised Hodge conjecture, in its $\ell$-adic form]\label{c:hodge}
For $X$ smooth projective, the coniveau filtration $\mathrm{Filt}^j H^i(\bar X, \mathbb Q_\ell)$, made of the classes that vanish on the complement of a closed subset of codimension $\ge j$, is the realisation of the largest sub-motive of $h^i(X)$ of level $\le i - 2j$.
\end{bigconj}

\src{\fac{12}, pp.~2--3, 41--42, 52--54, inventory dating 1964--[c.~1971]; \fac{36}, pp.~12, 16. The pages do not name the conjecture; the name is the reading's.}

\entry{The route.} Page~3 states the expected equality and draws its consequence, granting resolution of singularities and the standard conjectures \q{(conséquence de la résolution et des conjectures standard)}: the dimensions of the $\mathrm{Filt}^j H^i$ are independent of $\ell$, since traces of motivic endomorphisms are rational and independent of~$\ell$ (pp.~29--30). One inclusion comes from the coniveau spectral sequence, read at $E_1$: for $X$ smooth projective the classes supported in codimension $p$ are Gysin images, of level $\le n - 2p$ (p.~41). For the level itself, pp.~52--54 read weight and effectivity on the eigenvalues of Frobenius, and folder~36 applies this after reduction to a finite field: a piece of level~$1$ has eigenvalues divisible by $q^{n-1}$, so a $p$-adic unit eigenvalue shows that a piece is not of level~$1$ (p.~16).

\entry{Where it stops.} The reverse inclusion, that the largest sub-motive of low level lies in the coniveau filtration, is never argued, and the Frobenius characterisation of p.~54 is struck. In folder~36 he cannot show, even granting the standard conjectures, that the level of a motive can only drop under specialisation: \q{la triste vérité est que je ne sais pas} (p.~16).

\entry{Today.} Open. Grothendieck's 1969 paper corrected Hodge's general conjecture of 1950, which is false, into this form.

\begin{bigconj}[The period conjecture]\label{c:per}
Over $\overline{\mathbb Q}$, the functor that sends a motive to its de Rham realisation, its Betti realisation and the comparison isomorphism between them is fully faithful; equivalently, the periods of a motive generate a field of transcendence degree equal to the dimension of its motivic Galois group.
\end{bigconj}

\src{\fac{13}, \emph{Groupe de Galois motivique}, pp.~30--34, 47--51, dated 1965; \fac{18}, pp.~10--12, 41--50, inventory dating 1969--[c.~1972].}

\entry{The route.} Folder~13 starts from a negative observation: the algebraic data of a realisation do not say which Hodge data come from motives, and must be completed by conditions of the type \q{arg.\ de transcendance} (p.~31). Pages~33--34 prove an abstract lemma: if $g$ places the second lattice and $H$ is the group generated by the closure of the elements $ts^{-1}$, then $V \cap T'(V) = V^H$, so full faithfulness becomes $V^H = V^G$ for every representation~$V$. Page~47 asks the question without filtration or Hodge structure, reduces it at once to $\overline{\mathbb Q}$ and to $V$ simple, and settles rank one: \q{c'est une puissance de Tate, et ça marche}. Folder~18 turns it into transcendence: the field generated by a generic period point has transcendence degree $\dim G - 1$ for a group $G$ attached to the motive (p.~11), which \q{contient le th.\ de Schneider} for elliptic curves without complex multiplication; elliptic curves are the test case (p.~12), and pp.~46--50 sketch an adelic analogue.

\entry{Where it stops.} No non-Tate motive is treated. For a product of two non-isogenous elliptic curves without complex multiplication his bound gives at most $5$ against the expected $6$ (\fac{18}, p.~44). Over a field of finite type the functor is not fully faithful, and the page breaks off on \q{idiot !} before hoping for a weaker result (\fac{13}, p.~48).

\entry{Today.} Open. For an elliptic curve with complex multiplication the periods and quasi-periods have transcendence degree~$2$ (Chudnovsky); without, the expected degree~$4$ is open. A Lean statement of the elliptic case is in the Tau Ceti roadmap on elliptic periods.

\begin{bigconj}[The homotopy hypothesis]\label{c:hh}
$n$-groupoids, up to $n$-equivalence, are the same as homotopy $n$-types, and $\infty$-groupoids the same as homotopy types.
\end{bigconj}

\src{\fac{134-2}, pp.~42--46, postscript of a letter to Breen of 5 February 1975; pp.~48--51, letter of 17 February; pp.~145--149 and 167--170, French original of the letter of July 1975; \fac{134-1}, p.~6, letter to Deligne of 27 August 1974.}

\entry{The route.} To a Kan complex $K$ he attaches an $n$-category $C_n(K)$ of points, paths and homotopies, the top level taken up to homotopy, and notes that $\pi_i(K) \cong \pi_i(C_n K)$ for $i \le n$ almost by definition (pp.~42--43); the two localised categories should then be equivalent. He uses the hypothesis to choose the definitions: the yoga \q{pourrait constituer un fil d'Ariane} for defining non-strict $n$-categories, their functors and their equivalences (p.~43). He checks the dictionary on familiar constructions (p.~45): loops and $\operatorname{Hom}(x,x)$, mapping spaces and functor categories, homotopy fibres and fibre products, $K(\pi, n)$ and $n$-gerbes, truncation and classes of top arrows. The abelian case comes first: strict Picard $n$-categories should be complexes of length~$n$, and the correspondence between simplicial sets and $\infty$-groupoids a set-theoretic Dold--Puppe (pp.~146--147); for groups in $n$-categories he wants a \q{non-commutative chain complex}, and for $n = 1$ gives crossed modules with their invariant in $H^3$ (pp.~148--149). The functor from simplicial sets to $n$-groupoids is not negotiable: a theory without it \q{me semblerait canularesque} (p.~146). On the topos side, locally constant stacks on $X$ should be functors on $\Pi_n(X)$, and $\Pi_n$ recovered from them as in Galois theory (pp.~48--51, 167--170).

\entry{Where it stops.} Nothing is proved: the definition of non-strict $n$-categories is the missing piece, and he writes of the \q{inquiétant mirage} of $n$-limits (p.~169, reading uncertain). He cannot define the classifying topos of an $n$-gerbe for $n \ge 2$ (p.~45), and asks Breen for candidates for the non-abelian complexes. \emph{Pursuing Stacks} (1983), not transcribed in the edition, takes the programme up with coherators and test categories.

\entry{Today.} The answer depends on the model. It is false for strict $n$-groupoids from $n = 3$ (Simpson) and open for Grothendieck's globular $\infty$-groupoids in Maltsiniotis's formulation; an OpenAI preprint of 24 September 2026, not peer-reviewed, announces a proof for the Ara--Henry coherators \cite{huacicm}.

\begin{bigconj}[The section conjecture]\label{c:sec}
For a hyperbolic curve $U$ over a field $K$ of finite type over $\mathbb Q$, every section of $\pi_1(U) \to \Gal(\bar K/K)$ comes, up to conjugacy, from a rational point of $U$ or from a rational point at infinity.
\end{bigconj}

\src{\fac{149}, pp.~26--127, undated: Theorem~3 (pp.~31--33), the \q{Conjecture fondamentale} (p.~47), Conjectures A, B and B$'$ (pp.~55--62), relative curves (pp.~63--64, where \q{Nos conjectures impliquent donc qu'il y a exactement deux sortes de telles sections}), specialisation (pp.~69--98), the birational form (pp.~123--127); \fac{142}, pp.~79--81.}

\entry{The route.} Only injectivity is proved, and arithmetically: he first tried a geometric proof, \q{Finalement, je n'arrive pas à le prouver par voie géométrique} (p.~26). He reduces to a semi-abelian variety and the Kummer sequences for all~$n$; a point with trivial class is infinitely divisible, hence zero since the group of rational points is finitely generated (pp.~31--33), and one prime $\ell$ is not enough (p.~33). The section conjecture is the case \q{$V$ a point} of his conjecture on morphisms (p.~48). The rest of the folder reformulates the conjecture. Sections at a cusp are told from rational points by filling the hole (p.~54). Conjecture~B recasts the statement as a fixed point: the action of the Galois group on the universal cover completed by its cusps has \q{un pt fixe et un seul} (p.~58). Over a model the sections become lifts to $\pi_1(S)$, and a local criterion says that a rational point extends over a discrete valuation ring exactly when its section kills the inertia (p.~73); this compares sections over $K$ and over the residue fields. Pages~98--103 try to build counterexamples from bad reduction, without success.

\entry{Where it stops.} No route to surjectivity appears. Mordell's conjecture, which Faltings proved in 1983, is never used; the only arithmetic input is the Mordell--Weil theorem, for injectivity. The case of the thrice-punctured sphere at a cusp rests on a relation the page calls \q{hautement plausible} (p.~60). On reduction he writes that he has no idea for the moment (p.~102), and the explicit cocycle of a point is called fascinating work not yet begun (\fac{142}, p.~81).

\entry{Today.} Open. Injectivity is known, as are a birational $p$-adic form (Koenigsmann) and the real analogue.

\begin{bigconj}[Ihara's question: $\Gal(\overline{\mathbb Q}/\mathbb Q)$ and the Teichmüller tower]\label{c:gt}
The outer action of $\Gal(\overline{\mathbb Q}/\mathbb Q)$ on the profinite fundamental groups of the moduli spaces $M_{g,\nu}$ identifies it with the automorphisms of the tower that respect its structure; in Drinfeld's form, $\Gal(\overline{\mathbb Q}/\mathbb Q) \to \widehat{GT}$ is an isomorphism.
\end{bigconj}

\src{\fac{142}, pp.~12--16, 21--41, 86--100, inventory dating [from 1978]; \fac{143}, pp.~28--30; \fac{145}, pp.~4--37; \fac{147}, pp.~59, 103, 148--151; \fac{144}, p.~4; \fac{146}; \fac{67}, p.~49, January 1984.}

\entry{The route.} The target is stated in folder~143: \q{On aimerait} that $\pi_1(M_{g,\nu})$ map isomorphically onto the outer automorphisms of the profinite surface group that satisfy the loop condition, which would describe the Galois group through them (p.~29). Folder~142 reduces the action on any curve to the thrice-punctured sphere, by Belyi and a Lefschetz-type curve (pp.~86--88), and lists the necessary conditions for an automorphism of the profinite groupoid to come from Galois: it lifts to the covers, stabilises the kernels of morphisms, is compatible with morphisms between covers, and satisfies the loop condition, \q{sûrement la plus importante de toutes, et de très loin !} (p.~99). He conjectures them sufficient for one curve (pp.~94--95, Problem~\ref{p:gal}). Since the first conditions cannot be written down directly in the parameters (p.~98), he turns to the moduli spaces, with a conjecture tied to $M_{0,4}$ and $M_{0,5}$ that the page announces and that is nearly illegible (pp.~98--99). The low levels are computed: automorphisms commuting with $\mathfrak S_3$ as triples satisfying one relation (folder~142), pairs $(p, g)$ with a marginal \q{$= \chi(u)$} (folder~145), presentations in modular dimension $0$ and~$1$ (folder~146). Folder~147 then plans the discrete tower: a more or less complete understanding of the cases of modular dimension at most~$2$ \q{donnera la clef pour une compréhension complète dans le cas général} (p.~59), a programme to present all the groups and their maps by generators and relations (p.~103), and a gluing theorem at infinity over the strata of codimension $1$ and~$2$ (pp.~149--151, Problem~\ref{p:nc}).

\entry{Where it stops.} The sufficiency of the conditions, the $M_{0,4}$--$M_{0,5}$ conjecture, the programme of folder~147 and the gluing theorem are announced, not proved. Page~59 of folder~147 speaks of the discrete groups, in the analytic setting; reading it as the two-level principle for the Galois action is the edition's. The last step, that an automorphism satisfying the relations of the tower comes from Galois, is not attempted anywhere in the notes.

\entry{Today.} Open; the map is injective by Belyi's theorem. The two-level principle for the discrete tower was proved by Hatcher, Lochak and Schneps \cite{hls}.

\begin{bigconj}[The $p$-curvature conjecture]\label{c:pcurv}
Let $S$ be integral, of finite type and dominant over $\Spec \mathbb Z$, $X/S$ smooth, and $E$ a vector bundle on $X$ with an integrable connection. If the $p$-curvature of $E$ vanishes at almost all closed points of $S$, then $E$ has finite monodromy at the generic point.
\end{bigconj}

\src{\fac{93}, pp.~188--212, \emph{Connexions de Gauss--Manin et équations de Picard--Fuchs}, inventory dating [1972]--1981. No other folder has it.}

\entry{The route.} Page~194 notes, without proof, that finite monodromy forces vanishing $p$-curvature at almost all closed points, and asks: \q{La réciproque est-elle vraie ?} It then splits rank one: \q{La question précédente se décompose en deux}. First, is the class of $E$ in the Picard group torsion? If so, $E = \mathcal O_X$ with connection $d + \omega$, and the question becomes whether a closed form that is locally logarithmic at almost all primes has a multiple that is logarithmic. In characteristic~$0$ the target is explicit: the logarithmic forms are a group $\mathbb Z^r$ by Rosenlicht, so the integrable classes form $(\mathbb Q/\mathbb Z)^r$ (pp.~191--192). Pages~195--197 extend the question to torsors under group schemes, and for an abelian scheme rewrite it on the universal vector extension of the dual: a section that is divisible by $p$ at every closed point, is it of finite order? He reduces it to a local-global statement about connected subgroups, and the case of unipotent rank~$0$ is the one on which, \q{si ce résultat était acquis}, everything reduces to the semi-abelian case. Pages~198--212 compute the condition $\tilde C\omega = \tilde\omega$ on the Cartier operator and the Verschiebung of the Picard functor with connection; that these serve the question is our reading, not the page's.

\entry{Where it stops.} The converse and each of its reductions stay questions. There is no Gauss--Manin case, no Chebotarev argument and no reduction to curves in the notes.

\entry{Today.} Open. Katz proved it for Gauss--Manin connections (1972), and it holds for connections with solvable monodromy, rank one included (Chudnovsky for curves, André, Bost).

\section{How the list was made, and its limits}

The fourteen problems were chosen from the register of findings among the rows still open and outside what Lean reached, and with a statement precise enough to be proved or refuted. Rows whose only interest is a slip on the page, a dating, or a statement now known to be in print were left out, as were the conjectures already in the published record, which \S\ref{sec:big} treats from the side of the routes the notes take, and the companion paper from the side of their first appearance \cite{huacicm}.

Three limits apply to every entry. The transcriptions are first passes by a language model, and doubtful readings are marked but not all checked against the facsimile; a statement may rest on a misread word, and the entries say where we know this. The literature searches were made by the same means, with full texts where they were open and abstracts otherwise; an entry that says \q{not found} reports a search, not an absence. And the steps marked \q{ours} have not been checked by a specialist.

Corrections, references and proofs are welcome at the address above. The register of findings, the Lean files and the notices of the folders are on the project's site, \url{https://grothendieck.commutator.io/}.

\subsection*{Changes}
26 October 2026: the pseudoline dichotomy of \fac{80} proved (\S\ref{sec:pseudo}); the count of Kostant quadruples derived, and Problem~\ref{p:kostant} replaced by the question Grothendieck asked Kostant; Problem~\ref{p:ext2} narrowed after Breen's 1975 theorem; references corrected after a citation check. 29 October 2026: \S\ref{sec:big} added, on the routes the notes take to seven published conjectures.

\subsection*{Declaration}
This note was drafted with a language model (Claude Opus 5.5, Anthropic) from the edition's registers: the findings, the modernised readings and the Lean files. The transcriptions and the modernised readings it relies on were also made with language models (Claude Opus~5 and Opus~5.5). The author received no funding for this work.

\begin{thebibliography}{99}

\bibitem{andre} Y.~André, \emph{Une introduction aux motifs (motifs purs, motifs mixtes, périodes)}, Panoramas et Synthèses~17, SMF, 2004.
\bibitem{baues} H.-J.~Baues, \emph{Homotopy Type and Homology}, Oxford Mathematical Monographs, 1996.
\bibitem{bokowskipokora} J.~Bokowski, P.~Pokora, On the Sylvester--Gallai and the orchard problem for pseudoline arrangements, \emph{Period.\ Math.\ Hungar.} 77 (2018), 164--174, arXiv:1607.05864.
\bibitem{breen1975} L.~Breen, Un théorème d'annulation pour certains $\Ext^i$ de faisceaux abéliens, \emph{Ann.\ Sci.\ ÉNS} (4) 8 (1975), 339--352.
\bibitem{breen1999} L.~Breen, Monoidal categories and multiextensions, \emph{Compositio Math.} 117 (1999), 295--335, arXiv:math/9809104.
\bibitem{caramello} O.~Caramello, De Morgan classifying toposes, \emph{Adv.\ Math.} 222 (2009), 2117--2144, arXiv:0808.1519.
\bibitem{clausenmathew} D.~Clausen, A.~Mathew, Hyperdescent and étale K-theory, \emph{Invent.\ Math.} 225 (2021), 981--1076, arXiv:1905.06611.
\bibitem{collas} B.~Collas, Anabelian perspectives in Galois-Teichmüller theory, arXiv:2603.02848, 2026.
\bibitem{hodge3} P.~Deligne, Théorie de Hodge~III, \emph{Publ.\ Math.\ IHÉS} 44 (1974), 5--77.
\bibitem{durfee} A.~H.~Durfee, Mixed Hodge structures on punctured neighborhoods, \emph{Duke Math.\ J.} 50 (1983), 1017--1040.
\bibitem{boseeverettwismath} P.~Bose, H.~Everett, S.~Wismath, Properties of arrangement graphs, \emph{Internat.\ J.\ Comput.\ Geom.\ Appl.} 13 (2003), 447--462.
\bibitem{eppstein} D.~Eppstein, Drawing arrangement graphs in small grids, or how to play Planarity, \emph{J.\ Graph Algorithms Appl.} 18 (2014), 211--231, arXiv:1308.0066.
\bibitem{felsnergoodman} S.~Felsner, J.~E.~Goodman, Pseudoline arrangements, in \emph{Handbook of Discrete and Computational Geometry}, 3rd ed., ch.~5, CRC Press, 2017.
\bibitem{friedmanlaza} R.~Friedman, R.~Laza, The higher Du Bois and higher rational properties for isolated singularities, \emph{J.\ Algebraic Geom.} 33 (2024), 493--520, arXiv:2207.07566.
\bibitem{gabberorgogozo} O.~Gabber, F.~Orgogozo, Sur la $p$-dimension des corps, \emph{Invent.\ Math.} 174 (2008), 47--80, arXiv:0706.0671.
\bibitem{gabberzavyalov} O.~Gabber, B.~Zavyalov, Algebraization techniques and rigid-analytic Artin--Grothendieck vanishing, \emph{Ann.\ Sci.\ ÉNS} 59 (2026), 387--436, arXiv:2402.08741.
\bibitem{hls} A.~Hatcher, P.~Lochak, L.~Schneps, On the Teichmüller tower of mapping class groups, \emph{J.\ reine angew.\ Math.} 521 (2000), 1--24.
\bibitem{huacicm} M.~Hua, Formalising a mathematician's working notes: Lean proofs from the Grothendieck archive, preprint, 2026, \url{https://grothendieck.commutator.io/article/grothendieck-cicm.pdf}.
\bibitem{illusie2023} L.~Illusie, Revisiting deformations of truncated Barsotti--Tate groups, talk notes, Chicago, 6 March 2023, \url{https://www.imo.universite-paris-saclay.fr/~luc.illusie/Illusie-Chicago-2023-1.pdf}.
\bibitem{elephant} P.~T.~Johnstone, \emph{Sketches of an Elephant: A Topos Theory Compendium}, 2 vols, Oxford Logic Guides 43--44, Oxford, 2002.
\bibitem{kellyrottenberg} L.~M.~Kelly, R.~Rottenberg, Simple points in pseudoline arrangements, \emph{Pacific J.\ Math.} 40 (1972), 617--622.
\bibitem{kleiman1968} S.~L.~Kleiman, Algebraic cycles and the Weil conjectures, in \emph{Dix exposés sur la cohomologie des schémas}, Adv.\ Stud.\ Pure Math.~3, North-Holland, 1968, 359--386.
\bibitem{kockmoerdijk1992} A.~Kock, I.~Moerdijk, Every étendue comes from a local equivalence relation, \emph{J.\ Pure Appl.\ Algebra} 82 (1992), 155--174.
\bibitem{kockmoerdijk1996} A.~Kock, I.~Moerdijk, Spaces with local equivalence relations, and their monodromy, \emph{Topology Appl.} 72 (1996), 47--78.
\bibitem{kostant} B.~Kostant, The principal three-dimensional subgroup and the Betti numbers of a complex simple Lie group, \emph{Amer.\ J.\ Math.} 81 (1959), 973--1032.
\bibitem{mathew} A.~Mathew, On $K(1)$-local TR, \emph{Compositio Math.} 157 (2021), 1079--1119, arXiv:2005.08744.
\bibitem{milneadt} J.~S.~Milne, \emph{Arithmetic Duality Theorems}, 2nd ed., BookSurge, 2006.
\bibitem{moerdijk1993} I.~Moerdijk, Foliations, groupoids and Grothendieck étendues, \emph{Rev.\ Acad.\ Cienc.\ Zaragoza} (2) 48 (1993), 5--33.
\bibitem{mohar1995} B.~Mohar, Uniqueness and minimality of large face-width embeddings of graphs, \emph{Combinatorica} 15 (1995), 541--556.
\bibitem{mohar-robertson-vitray} B.~Mohar, N.~Robertson, R.~P.~Vitray, Planar graphs on the projective plane, \emph{Discrete Math.} 149 (1996), 141--157.
\bibitem{poonenstoll} B.~Poonen, M.~Stoll, The Cassels--Tate pairing on polarized abelian varieties, \emph{Ann.\ of Math.} 150 (1999), 1109--1149.
\bibitem{pop2019} F.~Pop, Finite tripod variants of I/OM: on Ihara's question/Oda--Matsumoto conjecture, \emph{Invent.\ Math.} 216 (2019), 745--797.
\bibitem{rosenthal} K.~I.~Rosenthal, Sheaves and local equivalence relations, \emph{Cahiers Top.\ Géom.\ Diff.} 25 (1984), 179--206.
\bibitem{saito} M.~Saito, Mixed Hodge modules, \emph{Publ.\ RIMS} 26 (1990), 221--333.
\bibitem{scheiderer} C.~Scheiderer, \emph{Real and Étale Cohomology}, Lecture Notes in Math.~1588, Springer, 1994.
\bibitem{sga4} M.~Artin, A.~Grothendieck, J.-L.~Verdier (eds.), \emph{Théorie des topos et cohomologie étale des schémas} (SGA~4), LNM 269, 270, 305, Springer, 1972--1973.
\bibitem{springer} T.~A.~Springer, Regular elements of finite reflection groups, \emph{Invent.\ Math.} 25 (1974), 159--198.
\bibitem{tonini2013} F.~Tonini, Stacks of ramified Galois covers, PhD thesis, 2013, arXiv:1307.1116.
\bibitem{tonini2017} F.~Tonini, Ramified Galois covers via monoidal functors, \emph{Transform.\ Groups} 22 (2017), 845--868, arXiv:1507.05309.
\bibitem{vitray} R.~P.~Vitray, The 2 and 3 representative projective planar embeddings, \emph{J.\ Combin.\ Theory Ser.\ B} 54 (1992), 1--12.

\end{thebibliography}

\end{document}
